\chapter{Sucesiones, límites y continuidad}
\label{chap:sucesiones-limites-continuidad}

El límite permite describir una tendencia sin confundirla con una igualdad en
un instante concreto. Empezaremos con sucesiones, donde el índice es discreto,
y trasladaremos después las mismas ideas a funciones de variable real. La
continuidad cerrará el capítulo al convertir información local sobre límites
en conclusiones globales de existencia.

\section{Sucesiones reales}
\label{sec:sucesiones-reales}

\subsection{Definición y formas de descripción}
\CanonicalUnit{CM-C02-U001}{}{}

\begin{definicion}[Sucesión real]
Una \emph{sucesión real} es una función
\[
  x:\N\longrightarrow\R, \qquad n\longmapsto x_n.
\]
El conjunto \(\N=\{1,2,\ldots\}\) es el conjunto de índices. La notación
\((x_n)_{n\ge1}\) representa la sucesión, mientras que \(x_n\) representa su
término de índice \(n\).
\end{definicion}

Una fórmula explícita determina cada término directamente; por ejemplo,
\(x_n=(-1)^n/n\). Una definición recursiva fija uno o varios términos
iniciales y expresa los siguientes a partir de los anteriores. La sucesión de
Fibonacci, por ejemplo, queda determinada por
\[
  F_1=F_2=1, \qquad F_{n+2}=F_{n+1}+F_n \quad(n\ge1).
\]
En cualquier definición recursiva hay que declarar los términos iniciales y
el conjunto de índices para el que la regla tiene sentido.

\subsection{Convergencia y divergencia}
\CanonicalUnit{CM-C02-U002}{}{}
\ProofRole{proof-006}{}

\begin{definicion}[Convergencia]
La sucesión \((x_n)\) \emph{converge} a \(L\in\R\) si, para todo
\(\varepsilon>0\), existe un índice umbral \(N\in\N\) tal que
\[
  n\ge N \quad\Longrightarrow\quad |x_n-L|<\varepsilon.
\]
Escribimos \(x_n\to L\) o \(\lim_{n\to\infty}x_n=L\). Si no existe tal
\(L\), la sucesión diverge.
\end{definicion}

El índice \(N\) puede depender de \(\varepsilon\), pero no del índice
cuantificado \(n\). La definición no exige que los primeros términos estén
cerca de \(L\); exige que todos los términos posteriores a un umbral lo estén.

\begin{ejemplo}
La sucesión \(x_n=1/n\) converge a cero. Dado \(\varepsilon>0\), elegimos un
entero \(N>1/\varepsilon\). Si \(n\ge N\), entonces
\[
  \left|\frac1n-0\right|=\frac1n\le\frac1N<\varepsilon.
\]
El testigo es \(N\); el símbolo \(n\) queda libre para representar cualquier
índice posterior.
\end{ejemplo}

La sucesión \((-1)^n\) no converge: sus términos pares valen \(1\) y los
impares valen \(-1\). Este ejemplo también muestra que una sucesión acotada no
tiene por qué ser convergente.

\subsection{Consecuencias y reglas de cálculo}
\CanonicalUnit{CM-C02-U003}{}{}
\ProofRole{proof-007}{}

\begin{teorema}[Unicidad del límite]
Una sucesión convergente tiene un único límite.
\end{teorema}

\begin{proof}
Supongamos que \(x_n\to L\) y \(x_n\to M\). Si \(L\ne M\), tomamos
\(\varepsilon=|L-M|/3\). Existen índices \(N_L\) y \(N_M\) que satisfacen
las respectivas definiciones. Para
\(n\ge N=\max\{N_L,N_M\}\), la desigualdad triangular da
\[
 |L-M|\le |L-x_n|+|x_n-M|<2\varepsilon
 =\frac23|L-M|,
\]
una contradicción. Por tanto, \(L=M\).
\end{proof}

\begin{proposicion}[Acotación]
Toda sucesión convergente es acotada.
\end{proposicion}

\begin{proof}
Si \(x_n\to L\), existe \(N\ge2\) tal que \(|x_n-L|<1\) para \(n\ge N\).
Entonces \(|x_n|<|L|+1\) en la cola. Los términos
\(x_1,\ldots,x_{N-1}\) forman un conjunto finito, de modo que
\[
 M=\max\{|x_1|,\ldots,|x_{N-1}|,|L|+1\}
\]
satisface \(|x_n|\le M\) para todo \(n\).
\end{proof}

\begin{teorema}[Álgebra y orden de límites de sucesiones]
Si \(x_n\to L\), \(y_n\to M\) y \(\lambda\in\R\), entonces
\[
 x_n+y_n\to L+M,\qquad
 \lambda x_n\to\lambda L,\qquad
 x_ny_n\to LM.
\]
Si \(M\ne0\), entonces \(y_n\ne0\) a partir de cierto índice y
\(x_n/y_n\to L/M\). Además, si \(x_n\le y_n\) a partir de cierto índice,
entonces \(L\le M\).
\end{teorema}

\begin{proof}[Demostración de las reglas algebraicas]
Para la suma se escribe
\[
 |(x_n+y_n)-(L+M)|
 \le |x_n-L|+|y_n-M|
\]
y se elige un único umbral posterior a los dos que corresponden a
\(\varepsilon/2\). La regla escalar se obtiene de
\(|\lambda x_n-\lambda L|=|\lambda|\,|x_n-L|\), con el caso
\(\lambda=0\) separado.

Para el producto, la convergencia de \((x_n)\) permite fijar \(K\ge1\) tal
que \(|x_n|\le K\). Entonces
\[
 |x_ny_n-LM|
 \le |x_n|\,|y_n-M|+|M|\,|x_n-L|.
\]
Las dos diferencias pueden hacerse menores que
\(\varepsilon/(2K)\) y \(\varepsilon/(2(|M|+1))\), respectivamente. Para el
cociente, si \(M\ne0\), a partir de cierto índice
\(|y_n|\ge|M|/2\), y por tanto
\[
 \left|\frac1{y_n}-\frac1M\right|
 =\frac{|y_n-M|}{|y_n|\,|M|}
 \le\frac{2}{|M|^2}|y_n-M|\longrightarrow0.
\]
La regla del producto aplicada a \(x_n(1/y_n)\) concluye la prueba.
\end{proof}

La misma elección de umbral debe servir simultáneamente para todas las
estimaciones que intervienen. Por ejemplo, al sumar dos convergencias se toma
\(N=\max\{N_1,N_2\}\), no un símbolo ya cuantificado. Las desigualdades
estrictas pueden perderse en el límite: \(0<1/n\) para todo \(n\), pero ambos
lados pueden tender a cero.

\begin{teorema}[Teorema del encaje para sucesiones]
Si \(a_n\le x_n\le b_n\) a partir de cierto índice y
\(a_n\to L\), \(b_n\to L\), entonces \(x_n\to L\).
\end{teorema}

\subsection{Monotonía, completitud y una sucesión iterativa}
\CanonicalUnit{CM-C02-U004}{}{}
\ProofRole{proof-008}{}
\ProofRole{proof-001}{}

Una sucesión es \emph{creciente} si \(x_{n+1}\ge x_n\) para todo \(n\), y
\emph{decreciente} si \(x_{n+1}\le x_n\). Diremos que es monótona si cumple
una de esas propiedades. Los términos \emph{estrictamente creciente} y
\emph{estrictamente decreciente} se reservan para las desigualdades estrictas.

\begin{teorema}[Convergencia monótona]
Toda sucesión creciente y acotada superiormente converge. Toda sucesión
decreciente y acotada inferiormente converge.
\end{teorema}

\begin{proof}[Estrategia de demostración]
Consideremos el caso creciente. Sea
\[
  S=\{x_n:n\in\N\}.
\]
Este conjunto es no vacío y está acotado superiormente. Por la
completitud de \(\R\), existe \(L=\sup S\). Dado \(\varepsilon>0\), el número
\(L-\varepsilon\) no es una cota superior de \(S\); por tanto existe \(N\)
tal que \(x_N>L-\varepsilon\). Como la sucesión es creciente, para \(n\ge N\)
se cumple
\[
  L-\varepsilon<x_N\le x_n\le L.
\]
Así, \(|x_n-L|<\varepsilon\), y \(x_n\to L\). El caso decreciente se reduce
al anterior aplicándolo a \((-x_n)\).
\end{proof}

El argumento muestra exactamente dónde interviene la completitud: convierte
la acotación en un número real \(L\) que puede actuar como límite.

\begin{ejemplo}[Una iteración]
Definimos
\[
  x_1=1, \qquad x_{n+1}=\sqrt{2+x_n}.
\]
Primero probamos que la sucesión está bien definida y acotada. Si
\(1\le x_n\le2\), entonces
\(1<\sqrt{2+x_n}\le2\); por inducción, \(1\le x_n\le2\) para todo \(n\).

Probamos ahora la monotonía. Como \(x_2=\sqrt3>x_1\), y la raíz cuadrada es
creciente, de \(x_n\ge x_{n-1}\) se deduce
\[
  x_{n+1}=\sqrt{2+x_n}\ge\sqrt{2+x_{n-1}}=x_n.
\]
La sucesión es creciente y está acotada superiormente; por el teorema anterior
converge a algún \(L\in[1,2]\). La sucesión desplazada \((x_{n+1})\) tiene el
mismo límite. Solo después de haber probado la convergencia usamos
\(x_{n+1}^2=2+x_n\) y las reglas algebraicas de límites para obtener
\[
  L^2=2+L, \qquad L^2-L-2=0.
\]
Las soluciones algebraicas son \(2\) y \(-1\), pero \(L\in[1,2]\), luego
\(L=2\).
\end{ejemplo}

\begin{advertencia}
Resolver la ecuación de punto fijo produce candidatos, no prueba que la
sucesión converja. La acotación y la monotonía son las que justifican la
existencia del límite en este ejemplo.
\end{advertencia}

\subsection{Sucesiones de Cauchy}
\CanonicalUnit{CM-C02-U005}{}{}
\ProofRole{proof-009}{}

\begin{definicion}[Sucesión de Cauchy]
Una sucesión \((x_n)\) es \emph{de Cauchy} si, para todo \(\varepsilon>0\),
existe \(N\in\N\) tal que
\[
  m,n\ge N \quad\Longrightarrow\quad |x_n-x_m|<\varepsilon.
\]
\end{definicion}

La definición compara términos de la sucesión entre sí y no presupone un
límite conocido. Para \(x_n=1/n\), si \(m,n\ge N\), entonces
\[
  \left|\frac1n-\frac1m\right|
  \le \frac1n+\frac1m
  \le \frac2N.
\]
Elegir \(N>2/\varepsilon\) prueba que la sucesión es de Cauchy sin introducir
ninguna igualdad de signo incorrecto.

Toda sucesión convergente es de Cauchy. En \(\R\) también vale el recíproco:

\begin{teorema}[Criterio de Cauchy en \(\R\)]
Una sucesión real converge si y solo si es de Cauchy.
\end{teorema}

La implicación difícil es una forma de completitud de \(\R\). No
desarrollaremos aquí su demostración técnica ni convertiremos el criterio en
un bloque independiente de ejercicios. Su función es reconocer que
``los términos se estabilizan entre sí'' basta para garantizar un límite real.

\subsection{Límites infinitos de sucesiones}
\CanonicalUnit{CM-C02-U006}{}{}

\begin{definicion}[Divergencia a infinito]
Escribimos \(x_n\to+\infty\) si, para todo \(M\in\R\), existe
\(N\in\N\) tal que
\[
  n\ge N \quad\Longrightarrow\quad x_n>M.
\]
Análogamente, \(x_n\to-\infty\) si para todo \(M\in\R\) existe \(N\) tal
que \(x_n<M\) cuando \(n\ge N\).
\end{definicion}

Los símbolos \(\pm\infty\) describen comportamiento límite; no son números
reales. Por eso las reglas que los contienen son abreviaturas de teoremas con
hipótesis de signo. Por ejemplo, si \(x_n\to+\infty\) y \(y_n\to L>0\),
entonces \(x_ny_n\to+\infty\). Si \(L<0\), el producto tiende a
\(-\infty\). Cuando \(L=0\), la expresión simbólica \(0\cdot\infty\) es una
forma indeterminada: no fija ningún resultado.

\subsection{Enriquecimiento: criterio de Stolz--Cesàro}
\CanonicalUnit{CM-C02-U007}{}{}
\ProofRole{proof-010}{}

El criterio siguiente no pertenece al núcleo de técnicas exigibles, pero
muestra cómo una razón de incrementos puede controlar una razón de sucesiones.

\begin{teorema}[Criterio de Stolz--Cesàro]
Sean \((a_n)\) y \((b_n)\) sucesiones reales. Supongamos que \((b_n)\) es
estrictamente creciente y \(b_n\to+\infty\). Si existe, finito o infinito, el
límite
\[
  \lim_{n\to\infty}
  \frac{a_{n+1}-a_n}{b_{n+1}-b_n}=L,
\]
entonces \(a_n/b_n\to L\).
\end{teorema}

Por ejemplo, con \(a_n=1+1/2+\cdots+1/n\) y \(b_n=n\), la razón de
incrementos es \(1/(n+1)\), que tiende a cero. El criterio da \(a_n/n\to0\).
Esta aplicación no usa continuidad del logaritmo ni anticipa resultados de
derivación.

\section{Límites de funciones}
\label{sec:limites-funciones}

\subsection{Puntos de acumulación y definición epsilon--delta}
\CanonicalUnit{CM-C02-U008}{}{}
\ProofRole{proof-012}{}

Sea \(D\subseteq\R\). Un punto \(a\in\R\) es un \emph{punto de acumulación}
de \(D\) si todo entorno reducido de \(a\) contiene algún punto de \(D\). No
se exige que \(a\in D\). Esta convención permite estudiar límites en un punto
eliminado y en un extremo del dominio.

\begin{definicion}[Límite relativo al dominio]
Sea \(f:D\to\R\) y sea \(a\) un punto de acumulación de \(D\). Decimos que
\(f(x)\) tiende a \(L\in\R\) cuando \(x\to a\), relativamente a \(D\), si
para todo \(\varepsilon>0\) existe \(\delta>0\) tal que, para \(x\in D\),
\[
  0<|x-a|<\delta
  \quad\Longrightarrow\quad
  |f(x)-L|<\varepsilon.
\]
Escribimos \(\lim_{x\to a,\,x\in D}f(x)=L\), omitiendo \(x\in D\) cuando el
dominio está claro.
\end{definicion}

El valor \(f(a)\), si existe, no interviene en el límite. Lo que importa es el
comportamiento de la función en puntos del dominio próximos pero distintos de
\(a\).

\begin{proposicion}[Unicidad del límite funcional]
Si el límite de \(f\) cuando \(x\to a\), relativamente al dominio, existe,
entonces es único.
\end{proposicion}

\begin{proof}
Supongamos que dos números distintos \(L\) y \(M\) satisfacen la definición y
tomemos \(\varepsilon=|L-M|/3\). Sean \(\delta_L\) y \(\delta_M\) los radios
correspondientes. Como \(a\) es punto de acumulación del dominio, existe
\(x\in D\) tal que
\[
  0<|x-a|<\min\{\delta_L,\delta_M\}.
\]
Entonces
\[
  |L-M|\le |L-f(x)|+|f(x)-M|<2\varepsilon,
\]
lo que contradice la elección de \(\varepsilon\).
\end{proof}

\begin{ejemplo}[Un argumento directo]
Para todo \(a\in\R\),
\(\lim_{x\to a}x^2=a^2\). Dado \(\varepsilon>0\), tomamos
\[
  \delta=\min\left\{1,\frac{\varepsilon}{2|a|+1}\right\}.
\]
Si \(|x-a|<\delta\le1\), entonces \(|x|\le|a|+1\), y
\[
 |x^2-a^2|=|x-a|\,|x+a|
 \le |x-a|(2|a|+1)<\varepsilon.
\]
La elección funciona también cuando \(a\le0\); no divide por \(a\).
\end{ejemplo}

\subsection{Caracterización mediante sucesiones}
\CanonicalUnit{CM-C02-U009}{}{}
\ProofRole{proof-011}{}

\begin{teorema}[Criterio secuencial]
Sea \(f:D\to\R\) y sea \(a\) un punto de acumulación de \(D\). Entonces
\[
 \lim_{x\to a,\,x\in D}f(x)=L
\]
si y solo si, para toda sucesión \((x_n)\) de puntos de
\(D\setminus\{a\}\) tal que \(x_n\to a\), se cumple \(f(x_n)\to L\).
\end{teorema}

\begin{proof}
Supongamos primero el límite funcional. Dado \(\varepsilon>0\), elegimos el
\(\delta\) de su definición. Como \(x_n\to a\), existe \(N\) tal que
\(|x_n-a|<\delta\) para \(n\ge N\); además, \(x_n\ne a\). Por tanto,
\(|f(x_n)-L|<\varepsilon\), y \(f(x_n)\to L\).

Para el recíproco, supongamos que el límite funcional no vale. Entonces existe
\(\varepsilon_0>0\) tal que, para cada \(n\in\N\), podemos elegir
\(x_n\in D\) con
\[
  0<|x_n-a|<\frac1n,
  \qquad
  |f(x_n)-L|\ge\varepsilon_0.
\]
Así, \(x_n\to a\), pero \(f(x_n)\) no converge a \(L\), contradiciendo la
hipótesis secuencial.
\end{proof}

El criterio permite demostrar que un límite no existe mediante dos sucesiones
que se acercan al mismo punto y producen límites distintos. Para
\(f(x)=\sin(1/x)\), \(x\ne0\), pueden elegirse índices positivos de forma que
\(1/x_n=\pi/2+2\pi n\) y \(1/y_n=3\pi/2+2\pi n\). Entonces
\(x_n,y_n\to0\), pero \(f(x_n)=1\) y \(f(y_n)=-1\).

\subsection{Límites laterales y acotación local}
\CanonicalUnit{CM-C02-U010}{}{}

El \emph{límite por la derecha} usa solamente puntos \(x>a\) del dominio:
\[
  \lim_{x\to a^+}f(x)=L.
\]
Formalmente es el límite relativo a \(D\cap(a,+\infty)\). El límite por la
izquierda, \(x\to a^-\), se define respecto de \(D\cap(-\infty,a)\). Cuando
ambos lados tienen puntos que se acumulan en \(a\), el límite bilateral existe
y vale \(L\) si y solo si los dos límites laterales existen y valen \(L\).

\begin{proposicion}[Acotación local]
Si \(\lim_{x\to a,\,x\in D}f(x)=L\in\R\), entonces existen
\(\delta>0\) y \(M>0\) tales que
\[
 x\in D,\quad 0<|x-a|<\delta
 \quad\Longrightarrow\quad |f(x)|\le M.
\]
Si además \(a\in D\), puede aumentarse \(M\) para incluir también \(f(a)\).
\end{proposicion}

\begin{proof}
Tomamos \(\varepsilon=1\). En un entorno reducido suficientemente pequeño,
\(|f(x)-L|<1\), luego \(|f(x)|<|L|+1\). Si \(a\in D\), basta reemplazar la
cota por \(\max\{|L|+1,|f(a)|\}\).
\end{proof}

\subsection{Reglas, orden y límite trigonométrico fundamental}
\CanonicalUnit{CM-C02-U011}{}{}
\ProofRole{proof-013}{}

Sean \(f,g:D\to\R\), y supongamos que, cuando \(x\to a\) relativamente a
\(D\),
\[
  f(x)\to L, \qquad g(x)\to M.
\]
Entonces se conservan la suma, el producto y los múltiplos escalares. Si
\(M\ne0\), entonces \(g(x)\ne0\) en algún entorno reducido relativo al
dominio y
\[
  \lim_{x\to a}\frac{f(x)}{g(x)}=\frac{L}{M}.
\]
La condición \(M\ne0\) es indispensable; que \(g(x)\ne0\) cerca de \(a\) no
basta por sí solo. Las reglas de orden y el teorema del encaje tienen la misma
forma que para sucesiones.

En concreto, si \(p(x)\le f(x)\le q(x)\) en algún entorno reducido relativo
al dominio y
\[
  \lim_{x\to a}p(x)=\lim_{x\to a}q(x)=L,
\]
entonces el teorema del encaje afirma que \(f(x)\to L\).

\begin{teorema}[Límite trigonométrico fundamental]
Si los ángulos se miden en radianes, entonces
\[
  \lim_{x\to0}\frac{\sin x}{x}=1.
\]
\end{teorema}

\begin{proof}
Para \(0<x<\pi/2\), la comparación geométrica de áreas en la circunferencia
unidad da
\[
  \sin x<x<\tan x.
\]
Al dividir por \(x>0\) y reorganizar,
\[
  \cos x<\frac{\sin x}{x}<1.
\]
Además,
\[
  0\le1-\cos x=2\sin^2(x/2)\le\frac{x^2}{2},
\]
donde se usa \(0\le\sin(x/2)\le x/2\). Por el teorema del encaje,
\(\cos x\to1\) y, de nuevo por encaje, \(\sin x/x\to1\) por la derecha.
Como \(\sin(-x)/(-x)=\sin x/x\), el cociente es par y el límite bilateral
tiene el mismo valor.
\end{proof}

Dos consecuencias que usaremos son
\[
  \sin x\sim x \quad(x\to0),
  \qquad
  1-\cos x=2\sin^2(x/2)=o(x) \quad(x\to0).
\]

\subsection{Límites infinitos, límites en el infinito y formas indeterminadas}
\CanonicalUnit{CM-C02-U012}{}{}
\ProofRole{proof-004}{}

Decimos que \(f(x)\to+\infty\) cuando \(x\to a\), relativamente a \(D\),
si para todo \(M\in\R\) existe \(\delta>0\) tal que
\[
 x\in D,\quad 0<|x-a|<\delta
 \quad\Longrightarrow\quad f(x)>M.
\]
Las definiciones de \(-\infty\) y de los límites laterales se obtienen
cambiando el signo o restringiendo el lado de aproximación.

Para límites en el infinito, por ejemplo,
\[
  \lim_{x\to+\infty}f(x)=L
\]
significa que para todo \(\varepsilon>0\) existe \(R>0\) tal que
\(x\in D\) y \(x>R\) implican \(|f(x)-L|<\varepsilon\). Se combinan de modo
análogo la aproximación \(x\to\pm\infty\) y los valores límite
\(L\), \(+\infty\) o \(-\infty\).

Expresiones como
\[
 \frac00,\quad \frac\infty\infty,\quad
 \infty-\infty,\quad 0\cdot\infty,
 \quad 1^\infty,\quad0^0,\quad\infty^0
\]
son \emph{formas indeterminadas}: describen información insuficiente sobre
los límites de los componentes, no operaciones entre números. Antes de aplicar
una técnica hay que identificar los límites separados y transformar la
expresión respetando su dominio real.

Para fijar una única caracterización principal del número \(e\), consideramos
\[
  u_n=\left(1+\frac1n\right)^n.
\]
La sucesión \((u_n)\) es creciente y está acotada superiormente; este resultado
se obtiene mediante desigualdades binomiales y se adopta aquí como parte del
límite de construcción de la exponencial real. Por completitud, \((u_n)\)
converge, y definimos
\[
  e=\lim_{n\to\infty}\left(1+\frac1n\right)^n.
\]
Las caracterizaciones mediante límites de variable real y derivadas se
deducirán cuando estén disponibles; no se tomarán como definiciones
independientes.

\subsection{Comparación asintótica y notación o pequeña}
\CanonicalUnit{CM-C02-U013}{}{}

Sea \(a\) un punto de acumulación del dominio común. Si \(g(x)\ne0\) en un
entorno reducido relativo al dominio, escribimos
\[
  f(x)=o(g(x)) \quad(x\to a)
\]
cuando \(f(x)/g(x)\to0\). Decimos que \(f\) y \(g\) son
\emph{asintóticamente equivalentes}, y escribimos \(f(x)\sim g(x)\), si
\[
  \frac{f(x)}{g(x)}\longrightarrow1.
\]
Si, además, ambas funciones tienden a cero, se llaman
\emph{infinitésimos equivalentes}. Un \emph{infinitésimo} es una función que
tiende a cero. Si
\(f=o(g)\), decimos que \(f\) es de orden superior a \(g\) en el punto y con
la variable de aproximación declarados.

Estas relaciones no son intercambiables. La afirmación \(f=o(g)\) es más
fuerte que una cota del tipo \(|f|\le C|g|\), y ninguna notación de comparación
se usará sin definición. Del límite trigonométrico fundamental se obtiene
\(\sin x\sim x\) cuando \(x\to0\); de
\(0\le1-\cos x\le x^2/2\) se deduce
\(1-\cos x=o(x)\).

\section{Continuidad}
\label{sec:continuidad}

\subsection{Continuidad en un punto y en un intervalo}
\CanonicalUnit{CM-C02-U014}{}{}

\begin{definicion}[Continuidad relativa al dominio]
Sea \(f:D\to\R\) y \(a\in D\). La función \(f\) es \emph{continua en
\(a\)} si, para todo \(\varepsilon>0\), existe \(\delta>0\) tal que
\[
 x\in D,\quad |x-a|<\delta
 \quad\Longrightarrow\quad |f(x)-f(a)|<\varepsilon.
\]
Si \(a\) es punto de acumulación de \(D\), esta condición equivale a
\[
  \lim_{\substack{x\to a\\x\in D}}f(x)=f(a).
\]
\end{definicion}

Una función es continua en un conjunto \(E\subseteq D\) si es continua en
cada punto de \(E\). En un intervalo cerrado \([a,b]\), la continuidad en
\(a\) se interpreta por la derecha y en \(b\) por la izquierda, siempre
relativamente al dominio. Denotaremos por \(C(D,\R)\) el espacio de funciones
continuas de \(D\) en \(\R\); dominio y codominio quedan así declarados.

\begin{teorema}[Álgebra y composición]
Si \(f,g:D\to\R\) son continuas en \(a\), también lo son
\(f+g\), \(\lambda f\) y \(fg\). Si \(g(a)\ne0\), el cociente \(f/g\) es
continuo en \(a\). Además, si \(f:A\to B\) es continua en \(a\),
\(g:C\to E\) es continua en \(f(a)\) y \(f(A)\subseteq C\), entonces
\(g\circ f\) es continua en \(a\).
\end{teorema}

La afirmación de composición respeta los tipos fijados en el
Capítulo~\ref{chap:numeros-reales-funciones}: primero se comprueba que los
valores de \(f\) pertenecen al dominio de \(g\), y después se componen los
límites.

\subsection{Continuidad de las funciones elementales}
\CanonicalUnit{CM-C02-U015}{}{}
\ProofRole{proof-014}{}

Las funciones constantes y la identidad son continuas. Las reglas algebraicas
implican que todo polinomio es continuo en \(\R\) y que una función racional
es continua en cada punto de su dominio natural. El valor absoluto es continuo
porque
\[
  \bigl||x|-|a|\bigr|\le|x-a|.
\]

Las funciones exponenciales, logarítmicas y trigonométricas son continuas en
sus dominios naturales. Como derivación representativa, para \(h\to0\),
\[
 \sin(a+h)-\sin a
 =\sin a(\cos h-1)+\cos a\sin h\longrightarrow0,
\]
por los límites trigonométricos anteriores. A partir de aquí usaremos la
continuidad de la familia elemental como herramienta, sin repetir una prueba
para cada fórmula. Las funciones inversas elementales son continuas en los
intervalos principales declarados en el Capítulo~1.

Una composición de funciones elementales es continua allí donde todas las
operaciones y funciones componentes estén definidas. La comprobación del
dominio precede siempre a la aplicación de este principio.

\subsection{Tipos de discontinuidad}
\CanonicalUnit{CM-C02-U016}{}{}

Sea \(f:D\to\R\) y \(a\) un punto relevante del dominio o de su frontera.
Usaremos la siguiente clasificación controlada:

\begin{itemize}
  \item Hay una \emph{discontinuidad evitable} si existe un límite finito
        \(L\) cuando \(x\to a\), pero \(f(a)\) no está definido o no vale
        \(L\). Redefinir \(f(a)=L\) elimina la discontinuidad.
  \item Hay una \emph{discontinuidad de salto} si los límites laterales
        finitos existen y son distintos. El tamaño del salto es la diferencia
        absoluta entre ellos.
  \item Hay una \emph{discontinuidad infinita} si al menos uno de los límites
        laterales es \(+\infty\) o \(-\infty\).
  \item Hay una \emph{discontinuidad oscilatoria} si no existe límite por una
        oscilación que persiste a toda escala.
\end{itemize}

La función \((x^2-1)/(x-1)\) en \(\R\setminus\{1\}\) tiene una
discontinuidad evitable en \(1\). La función signo tiene un salto en \(0\),
\(1/x^2\) tiene una discontinuidad infinita allí y \(\sin(1/x)\) presenta una
discontinuidad oscilatoria. La \emph{función parte entera inferior}, denotada
\(\lfloor x\rfloor\), tiene saltos en los enteros; no usaremos \([x]\) para
representarla.

Para funciones monótonas definidas en un intervalo, los límites laterales
finitos existen en los puntos interiores; las discontinuidades que aparecen
allí son de salto. La hipótesis de dominio intervalo no puede omitirse en favor
de un dominio arbitrario.

\section{Teoremas de existencia en intervalos cerrados}
\label{sec:teoremas-existencia}

\subsection{Teorema del valor extremo}
\CanonicalUnit{CM-C02-U017}{}{}
\ProofRole{proof-015}{}

\begin{teorema}[Teorema del valor extremo de Weierstrass]
Si \(f:[a,b]\to\R\) es continua y \(a<b\), entonces existen
\(x_{\min},x_{\max}\in[a,b]\) tales que
\[
  f(x_{\min})\le f(x)\le f(x_{\max})
  \qquad\text{para todo }x\in[a,b].
\]
Es decir, \(f\) alcanza un mínimo absoluto y un máximo absoluto.
\end{teorema}

El teorema combina continuidad con un intervalo cerrado y acotado. Ninguna de
estas condiciones es decorativa. En \((0,1)\), la función \(f(x)=x\) no
alcanza máximo ni mínimo. En \([0,1]\), la función definida por
\(f(0)=0\) y \(f(x)=1/x\) para \(x>0\) no es continua y tampoco está
acotada. El resultado garantiza existencia; no proporciona por sí solo un
método para localizar los extremos.

\subsection{Teorema de Bolzano}
\CanonicalUnit{CM-C02-U018}{}{}
\ProofRole{proof-016}{}
\ProofRole{proof-001}{}

\begin{teorema}[Teorema de Bolzano]
Sea \(f:[a,b]\to\R\) continua, con \(a<b\). Si
\(f(a)f(b)<0\), entonces existe \(c\in(a,b)\) tal que \(f(c)=0\).
\end{teorema}

El teorema afirma existencia, no unicidad. Por ejemplo, una función continua
puede cambiar de signo entre los extremos y tener varios ceros interiores.

\begin{proof}[Estrategia de bisección y completitud]
Se divide \([a,b]\) por su punto medio. Si el valor en ese punto es cero, la
conclusión ya está probada; en caso contrario, se conserva una mitad cuyos
valores extremos tienen signos opuestos. Al repetir el proceso se obtienen
intervalos cerrados encajados
\([a_n,b_n]\), con longitudes \((b-a)/2^n\), y con
\(f(a_n)f(b_n)<0\). La completitud de \(\R\) garantiza un único punto
\(c\) común a todos ellos; equivalentemente, \(a_n\to c\) y \(b_n\to c\).
Por continuidad, \(f(a_n)\to f(c)\) y \(f(b_n)\to f(c)\). Como los valores
conservan signos opuestos, el único límite común posible es \(f(c)=0\).
\end{proof}

La bisección es también un puente natural hacia el cálculo numérico: produce
intervalos que encierran un cero y un control explícito de su anchura. El
diseño de un algoritmo y sus actividades computacionales pertenece a las
prácticas posteriores, no a este capítulo compartido.

Como aplicación, todo polinomio real de grado impar tiene al menos una raíz
real. Si su coeficiente principal es positivo, el polinomio tiende a
\(+\infty\) cuando \(x\to+\infty\) y a \(-\infty\) cuando
\(x\to-\infty\); para coeficiente principal negativo se invierten los signos.
En ambos casos existen extremos finitos con valores de signo contrario, y
Bolzano proporciona un cero entre ellos.

\subsection{Teorema del valor intermedio}
\CanonicalUnit{CM-C02-U019}{}{}
\ProofRole{proof-017}{}

\begin{teorema}[Teorema del valor intermedio]
Sea \(f:[a,b]\to\R\) continua. Si \(z\) está entre \(f(a)\) y \(f(b)\),
entonces existe \(c\in[a,b]\) tal que \(f(c)=z\).
\end{teorema}

\begin{proof}
Si \(z=f(a)\) o \(z=f(b)\), basta elegir el extremo correspondiente.
Supongamos que \(z\) está estrictamente entre ambos valores y definamos
\[
  g(x)=f(x)-z.
\]
La función \(g\) es continua en \([a,b]\), y \(g(a)\) y \(g(b)\) tienen
signos opuestos. Por el teorema de Bolzano, existe \(c\in(a,b)\) tal que
\(g(c)=0\). Por tanto, \(f(c)=z\).
\end{proof}

La continuidad local en cada punto del intervalo es la hipótesis que permite
obtener esta conclusión global. El teorema tampoco afirma unicidad: una misma
altura puede alcanzarse en varios puntos.
